\vspace{-2mm}
\section{Limitations and Future Work}\label{sec:limit}
\vspace{-1mm}
While our results demonstrate the effectiveness of \toolName{}, several directions remain open for future exploration. First, admission-time verification trades a modest amount of overhead for stronger reliability guarantees; lighter-weight verifiers, including learned models or rule-based checks for common claim types, could further improve efficiency while preserving the grounding benefits we observe. Second, \toolName{} inherits the decomposition quality of agents. Overly coarse decompositions leave agents with under-specified subtasks, while overly aggressive decompositions can spawn unnecessary agents and overcomplicate reasoning. A promising direction is to train adaptive agents that decide when to split, merge, or terminate subtasks based on the shared context. Third, as noted in prior work~\citep{zhang2025recursive}, there is no universally optimal prompt across models; different model families may require tailored prompts to elicit the intended behavior. Combining \toolName{} with prompt-evolution methods such as GEPA~\citep{agrawal2025gepa} could further adapt its decomposition, summarization, and verification prompts to each model family.

Beyond the benchmarks studied here, \toolName{} suggests a promising direction for automated research systems~\citep{alphaevolve2025, liu2026skydiscover}. Research workflows naturally combine test-time scaling with long-context reasoning: agents must explore alternative hypotheses, inspect large collections of papers or experimental logs, compare evidence across sources, and iteratively refine conclusions. A decentralized shared context could make such systems more efficient by preventing agents from repeatedly reading the same papers or rerunning the same failed analyses; more effective by allowing useful findings to propagate across parallel research threads; and more robust by admitting only verified claims into the shared state. We therefore view automated research as a natural application domain for future decentralized MAS.

\vspace{-2mm}
\section{Conclusion}
\vspace{-1mm}
We introduce \toolName{}, a decentralized multi-agent system in which agents coordinate through a shared context and task queue rather than a central controller. By admitting only compact, verified updates into shared state, \toolName{} turns intermediate progress into reusable problem state: agents can build on prior findings, avoid repeated failures, preserve constraints, and recover detailed evidence only when needed. On SWE-bench Verified, \toolName{} improves test-time scaling by making discoveries and failures reusable across parallel attempts, achieving higher pass rates while reducing cost. On LongBench-v2 Multi-Doc QA, \toolName{} improves long-context reasoning by constructing a verified, hierarchical view of the corpus, leading to higher accuracy. We further show that \toolName{} is complementary to RLM, where decentralized validated state improves programmatic aggregation. These results suggest that scalable multi-agent systems require not only more parallel agents, but also a reliable communication substrate for sharing progress across them.


\section{Acknowledgments}
We thank Shayan Talaei, Jon Saad-Falcon, Jacky Kwok, Hermann Kumbong, Ishan Khare, Miria Feng, Qizheng Zhang, Swapnil Gandhi, Michael Y. Li, Chun Deng, Chong Zeng, Rui Li, Ke Li, Marquita Ellis, Hangoo Kang, Adrian Gamarra Lafuente, Charles Ding, Tarun Suresh, Ziyu Chen, Zhuohan Gu, Yize Liu, and Ligeng Zhu. We would also like to thank our collaborators at the Stanford Artificial Intelligence Laboratory (SAIL) and Stanford HAI.

We gratefully acknowledge support from federal sources: NSF under No. 24-554 (AIMing) and DARPA under No. HR00112520038 (Fallingwater). We also gratefully acknowledge support from Stanford HAI, IBM (Code Generation), Lightspeed, Google, and Google DeepMind.